%\begin{mdframed}
%    \textbf{La extensión máxima para esta sección es un párrafo breve (máximo 1/4 de página).}
%\end{mdframed}

%Escriba una conclusión breve que interprete los resultados más importantes. Cada afirmación debe estar respaldada por los datos presentados.


%Los greedy eran casi siempre más rapidos que el fuerza bruta, pero a medida que aumentaba la cantidad de animes el fuerza bruta crecio rapido



Respecto a los resultados obtenidos, el gráfico de tiempo deja claro que Fuerza Bruta crece demasiado rápido, lo cual lo hace inviable para casos grandes, a diferencia de los algoritmos Greedy que son muy veloces y escalan sin problemas. En cuanto a la satisfacción, era de esperarse que Fuerza Bruta consiga siempre el puntaje máximo. Sin embargo, aunque los Greedy no llegan a ese óptimo, logran resultados bastante cercanos en una fracción del tiempo. Finalmente, al comparar ambos algoritmos Greedy, puede notarse claramente en el último gráfico que \texttt{Greedy\_Menos\_Capitulos} es el ganador, ya que logra una mayor satisfacción que el otro Greedy en casi todos los tamaños probados.
\\
Respecto al espacio, no se mostró acá ya que en todas las mediciones dio 0, esto es debido al trabajo interno del sistema operativo, que se encarga de asignar el espacio. Este programa no llega a completar el espacio que se le asignó de base, así que nunca pide memoria extra lo cual dificulta mucho la medición del espacio utilizado.












